\section{Panorama de la industria: quién compite por el ciclo cerrado de datos de la robótica y la IA física}

En la entrevista se mencionan varios grupos de actores: los equipos de modelos grandes, los equipos de modelos del mundo, los equipos de VLA, los fabricantes de robots y las empresas de simulación y datos sintéticos. No forman una simple cadena de proveedores y clientes, sino una relación de simbiosis. Quien controle el ciclo cerrado de datos tiene más probabilidades de controlar el cerebro del robot.

\begin{figure}[H]
\centering
\includegraphics[width=0.95\textwidth]{data-industry-landscape.png}
\caption{Panorama de la industria de datos: modelos grandes, modelos del mundo, VLA, fabricantes de robots y empresas de datos de simulación.}
\end{figure}

\begin{knowledgebox}{Lectura de la figura: por qué es una red simbiótica}
Los equipos de modelos grandes aportan el modelo base y la infraestructura; los equipos de modelos del mundo se ocupan de la predicción física; los equipos de VLA conectan con las acciones; los fabricantes de robots tienen los robots reales; las empresas de simulación aportan escenarios a escala. Ninguna de las partes puede completar por sí sola todo el ciclo cerrado de datos.
\end{knowledgebox}

\subsection{Fabricantes de robots frente a empresas del cerebro robótico}

En la conducción autónoma, las empresas que tienen una flota de vehículos cuentan de forma natural con datos del lado del dispositivo; en la robótica todavía no existe un despliegue a gran escala en dispositivos, por lo que los fabricantes de robots no necesariamente tienen una ventaja de datos por defecto. Si una empresa del cerebro robótico controla la simulación y la infraestructura, también puede desarrollar capacidades primero; si un fabricante de robots logra formar un ciclo cerrado con datos reales, también será muy relevante.

\begin{warningbox}{La experiencia de la conducción autónoma no se traslada directamente a la robótica}
La conducción autónoma cuenta con entornos viales relativamente estandarizados y con flotas de gran escala; en la robótica, los cuerpos de los robots, las tareas y los escenarios están mucho más dispersos. El motor de datos debe rediseñarse; no se puede copiar la ruta de FSD.
\end{warningbox}

\subsection{Resumen de la sección}

La industria de datos para robótica será una red simbiótica. Los modelos, los modelos del mundo, los VLA, los fabricantes de robots y las empresas de datos de simulación pueden convertirse en nodos clave; lo decisivo es quién logra formar un ciclo cerrado.
